\documentclass[a4paper,10pt]{report}\input{../0.Préambule/Préambule_cours_prof}\begin{document}

\newpage
\section*{Les parallélogrammes}
\addcontentsline{toc}{section}{Les parallélogrammes}

\subsection*{Propriétés des parallélogrammes}
\addcontentsline{toc}{subsection}{Propriétés des parallélogrammes}

\begin{exer1}

\begin{flushright}
\includegraphics[scale=1]{../40.Image/15. Exo-Les parallélogrammes 1.png}
\end{flushright}

\vspace{-2.5cm}
\begin{enumerate}[font=\color{exer1}]
\item Écris tous les noms possibles du parallélogramme ci-contre.
\item Sur la figure, repasse en vert le coté opposé à $[PA]$, en bleu un coté \\ consécutif à $[PA]$, en rouge l'angle opposé à $\widehat{PCR}$ et en violet un  \\ angle consécutif à $\widehat{PAR}$.
\item Complète les phrases suivantes avec l'un ou plusieurs des mots : \begin{center}
consécutifs \quad \quad diagonales \quad \quad opposés \quad \quad cotés  \quad \quad  angles.
\end{center}

Dans le parallélogramme $PARC$ :
\begin{enumerate}[font=\color{exer1}]
\item Les \dotfill $[PA]$ et $[CR]$ sont de même mesure.
\item La somme des deux \dotfill $\widehat{PAR}$ et $\widehat{ARC}$ fait $180$\degre .
\item Les \dotfill $[PR]$ et $[CA]$ se coupent en leur milieu.
\item Les \dotfill $\widehat{CPA}$ et $\widehat{CRA}$ sont de même mesure.
\item Les \dotfill $[PC]$ et $[RA]$ sont parallèles.
\end{enumerate}
\end{enumerate}
\end{exer1}

\medskip
\begin{exer1}
Dans la figure ci dessous, les droites d'une \\ même couleur sont parallèles.

\vspace{-1.2cm}
\begin{flushright}
\includegraphics[scale=1]{../40.Image/15. Exo-Les parallélogrammes 2.png}
\end{flushright}

\vspace{-2.8cm}
\begin{enumerate}[font=\color{exer1}]
\item Nomme tous les parallélogrammes de cette figure.
\item Pourquoi peux-tu affirmer que ce sont des \\ parallélogrammes ?
\end{enumerate}
\end{exer1}

\vspace{0.5cm}
\begin{exer1}
Dessine et code le parallélogramme $ABCD$ selon les consignes et cite dans chaque cas la propriété du parallélogramme qui t'a permis de réaliser la consigne.

\begin{center}
\renewcommand{\arraystretch}{1.8}
\begin{tabularx}{\linewidth}{|>{\centering}m{4cm}|>{\centering}m{4.5cm}|>{\centering \arraybackslash}X|}
\hline
\rowcolor{exer1!20}Figure & Consigne & Propriété\\\hline
 & Code les cotés de même longueur & \\\cline{2-3}
 & Colorie d'une même couleur les angles de même mesure & \\\cline{2-3}
 & Code les longueurs égales sur les diagonales & \\\hline
\end{tabularx}
\end{center}
\end{exer1}

\noindent \begin{minipage}[t]{8cm}
\begin{exer1}
Complète les étiquettes sachant que $ROSE$ est un parallélogramme. Justifie tes réponses.

\begin{center}
\includegraphics[scale=1.2]{../40.Image/15. Exo-Les parallélogrammes 3.png}
\end{center}
\end{exer1}
\end{minipage}
\hspace{5mm}
\begin{minipage}[t]{8.5cm}
\begin{exer1}
La figure est dessiné a main levée

\begin{center}
\includegraphics[scale=0.8]{../40.Image/15. Exo-Les parallélogrammes 4.png}
\end{center}

\begin{enumerate}[font=\color{exer1}]
\item Complète les étiquettes sachant que $BLEU$ est un parallélogramme.
\item Justifie ta réponse pour l'angle $\widehat{BLE}$
\item Justifie ta réponse pour la longueur $BU$
\end{enumerate}
\end{exer1}
\end{minipage}
 
\begin{exer2}
On considère le parallélogramme $ABCD$

\vspace{-3mm}
\noindent\begin{minipage}{10cm}
\begin{enumerate}[font=\color{exer2}]
\item Quel est la mesure de l'angle $\widehat{CBA}$ ?
\item Pourquoi ?
\end{enumerate}
\end{minipage}
\hspace{5mm}
\begin{minipage}{6cm}
\includegraphics[scale=1]{../40.Image/15. Exo-Les parallélogrammes 5.png} 
\end{minipage}
\end{exer2}

\begin{exer2}
On considère le parallélogramme $UVWT$

\vspace{-3mm}
\noindent\begin{minipage}{10cm}
\begin{enumerate}[font=\color{exer2}]
\item Quel est la longueur $TW$ ?
\item Pourquoi ?
\end{enumerate}
\end{minipage}
\hspace{5mm}
\begin{minipage}{6cm}
\includegraphics[scale=1]{../40.Image/15. Exo-Les parallélogrammes 6.png}
\end{minipage}
\end{exer2}

\vspace{-3mm}
\begin{exer2}
$EFGH$ est un parallélogramme.

\begin{enumerate}[font=\color{exer2}]
\item Justifie que $(EF)$ est parallèle à $(GH)$ et que $(EH)$ est parallèle à $(FG)$
\item Quelle est la somme des quatre angles de ce parallélogramme ? Justifie.
\end{enumerate}
\end{exer2}

\begin{exer3}
$ABCD$ est un parallélogramme de centre $O$. Justifie que $O$ est le milieu du segment $[AD]$
\end{exer3}

\subsection*{Construire des parallélogrammes}
\addcontentsline{toc}{subsection}{Construire des parallélogrammes}

\begin{exer1}
Sur la figure ci-contre, trace à main levée :

\vspace{-3mm}
\noindent \begin{minipage}{10cm}
\begin{enumerate}[font=\color{exer1}]
\item en bleu, le point $D$ tel que $ABCD$ soit un parallélogramme.
\item en vert, le point $E$ tel que $AEBC$ soit un parallélogramme.
\item en rouge, le point $F$ tel que $ABFC$ soit un parallélogramme.
\end{enumerate}
\end{minipage}
\hspace{25mm}
\begin{minipage}{6cm}
\includegraphics[scale=0.2]{../40.Image/15. Exo-Les parallélogrammes 7.png} 
\end{minipage}
\end{exer1}

\vspace{-3mm}
\begin{exer1}
Place les points $H$ et $K$, pour que $EFHG$ et $JKL$ soient des parallélogrammes :

\begin{center}
 \includegraphics[scale=0.4]{../40.Image/15. Exo-Les parallélogrammes 8.png}
 \end{center} 
\end{exer1}

\begin{exer1}
Avec l'équerre et la règle non graduée, place dans chaque cas le point $P$ pour que $MNOP$ soit un parallélogramme.

\vspace{-4mm}
\begin{center}
\includegraphics[scale=0.3]{../40.Image/15. Exo-Les parallélogrammes 9.png} 
\end{center}
\end{exer1}

\begin{exer1}
Avec le compas, place dans chaque cas le point $S$ pour que $RSTU$ soit un parallélogramme.

\begin{center}
\includegraphics[scale=0.3]{../40.Image/15. Exo-Les parallélogrammes 10.png} 
\end{center}
\end{exer1}

%\begin{exer2}
%Dans chaque cas, place les points $M$ et $N$, pour que $KLMN$ soit un parallélogramme de centre $O$ :

%\includegraphics[scale=0.4]{../40.Image/15. Exo-Les parallélogrammes 11.png} 
%\end{exer2}
\end{document}